\chapter{Filter Choice}
\label{ch:bf-filter-choice}

Filter choice is the capstone decision register for the preceding particle,
transport, HMC-target, and verification staircase.  It is not only a speed
decision.  For HMC, the selected backend determines the log likelihood, the
derivative path, the finite-failure policy, and the compilation surface.  One
should choose the simplest backend that gives the required target fidelity and a
validated gradient.

\section{Decision Register}
\label{sec:bf-filter-choice-register}

The decision should be recorded in a short register:
\begin{description}
  \item[Exact Kalman, solve, or square-root LGSSM]
  Best for linear Gaussian likelihoods and regression oracles.  The main risks
  are SPD failures, mask handling, and large-scale factor failures.  The
  relevant check is analytic/custom score parity plus compiled value-gradient
  parity.

  \item[EKF]
  Best as a cheap nonlinear baseline and affine recovery check.  The main risks
  are linearization bias and Jacobian fragility.  The relevant check is affine
  Kalman recovery plus finite-gradient stress tests.

  \item[Sigma-point]
  Best for moderate nonlinear Gaussian approximations.  The main risks are an
  approximate target and covariance-update instability.  The relevant check is
  reference recovery, static sigma-point shapes, and derivative parity.

  \item[Square-root sigma-point]
  Best for value-side robustness under near-singular covariance.  The main risk
  is false confidence when the implementation reconstructs covariance matrices
  internally.  The relevant check is factor reconstruction, finite gradients,
  and compiled parity.

  \item[SVD sigma-point]
  Historical eigenderivative UKF/SVD routes are now retained primarily for
  diagnostic comparison, regression baselines, and bounded small-case oracles.
  The main risks are close or repeated spectral values and raw autodiff
  instability.  They should not be the preferred HMC-facing route once the
  strict-SPD principal-square-root backend is available.  The remaining role of
  the historical branch is eigen-gap telemetry, failure-mode diagnosis, and
  comparison against the promoted square-root branch.

  \item[Particle or differentiable particle filter]
  Best for non-Gaussian diagnostics and frontier approximations.  The main
  risks are Monte Carlo variance, resampling non-differentiability, and relaxed
  target bias.  The relevant check is an estimator-variance report together with
  a clearly declared correction policy.
\end{description}

\section{Recommended Order}
\label{sec:bf-filter-choice-order}

For a new model, BayesFilter should proceed in this order:
\begin{enumerate}[label=(\roman*)]
  \item reduce to an LGSSM or affine test case and validate the exact Kalman
    likelihood and derivatives;
  \item add masks, missing-data patterns, and mixed-frequency timing while
    retaining exact references;
  \item test EKF or sigma-point approximations on low-dimensional nonlinear
    examples with dense or trusted references;
  \item only then attempt DSGE-scale SVD sigma-point HMC;
  \item treat particle methods and transport methods as separate evidence
    streams unless a correction policy makes them target-preserving.
\end{enumerate}

\section{Industrial Default}
\label{sec:bf-filter-choice-industrial-default}

For large industrial applications, the default should be an exact or controlled
approximate likelihood with an analytic or custom derivative path.  Raw tape
gradients through spectral decompositions are acceptable as diagnostic tools and
small-case oracles, but they should not be the final route unless spectral-gap,
finite-gradient, and compiled parity evidence survives the intended stress
regime.

This policy is intentionally conservative.  The SVD sigma-point debugging work
showed that value robustness and HMC-gradient robustness are different
problems.  The next industrial and case-study chapters use this register as the
entry gate for application-specific backend choice rather than discovering those
distinctions only after long chains.
